immini[thm]{\subsection*{I. Các thông số trạng thái của một lượng khí}
\newghichu{
$\bullet$ Một lượng khí đựng trong một bình kín được xác định bởi bốn đại lượng là
\begin{itemize}
	\item khối lượng m [kg]
	\item thể tích V [$m^3$]
	\item nhiệt độ T [K]
	\item áp suất p [$Pa$]
\end{itemize}
$\bullet$ $(V,\ T,\ p)$ là các thông số trạng thái của một khối lượng khí xác định.
}}{includegraphics[width=4cm]{Images/KNTT91T}}
\Binhluan{\begin{itemize}
	\item Khí chuyển từ trạng thái này sang trạng thái khác bằng các quá trình biến đổi trạng thái, gọi tắt là các quá trình.
	\begin{center}
		includegraphics[width=8cm]{Images/KN92T}
	\end{center}\end{itemize}}{Áp suất chất khí được hiểu là áp suất trung bình của các phân tử khí tác dụng lên thành bình}

\begin{itemize}	\item Trong hầu hết các quá trình biến đổi trạng thái của một khối lượng khí xác định thì cả ba thông số $(p,\ V,\ T)$ đều có thể biến đổi.
	\item Những quá trình đơn giản mà chỉ có hai thông số biến đổi còn một thông số không đổi gọi là các \textit{đẳng quá trình}.
\end{itemize}

\begin{note}
\textbf{Áp suất} là độ lớn của áp lực trên một đơn vị diện tích bị ép.
\begin{align*}
	\boder{p=\dfrac{F}{S}}\quad \ct{trong đó}\ 
	\begin{cases}
		F:\ \ct{áp lực}\quad (N)\\
		S:\ \ct{diện tích mặt bị ép}\quad (m^2)\\
		p:\ \ct{áp suất}\quad (N/m^2\ \ct{hay}\ Pa)
	\end{cases}
\end{align*}
\end{note}
\begin{bang}[tabularx=l|l|X|X|X|X]{Những đơn vị đổi trong chất khí}
	&\textbf{SI}&Công thức&&&\\\hline
	\textbf{p}&$Pa$&&$1\ Pa=1\ N/m^2$&$1\ atm=1,03.10^5\ Pa$&$1\ mmHg=133,3\ Pa$\\\hline
	\textbf{T}&$K$&$T=t+273$&$0^\circ C=273\ K$&$0\ K=-273^\circ C$&\\\hline
	\textbf{V}&$m^3$&&$1\ l=1\ dm^3$&$1\ ml=1\ cm^3$&$1\ cc=1\ cm^3$
\end{bang}
\newpage 
immini[thm]{\subsection*{II. Định luật Boyle}
\subsubsection{Quá trình đẳng nhiệt}
\newghichu{
Quá trình đẳng nhiệt là quá trình biến đổi trạng thái của chất khí trong đó nhiệt độ được giữ không đổi.
}
\subsubsection{Định luật Boyle}
\begin{itemize}
	\item Trong quá trình đẳng nhiệt của một khối lượng khí xác định, áp suất tỉ lệ nghịch với thể tích
	\begin{center}
		\boder{
			p\sim \dfrac{1}{V}\ \text{hay}\ pV = \ct{hằng số}
		}
	\end{center}
	\item Vậy, khi nhiệt độ khối khí không đổi thì ta có hệ quả:
	\begin{center}
		\boder{
			p_{1}V_{1} = p_{2}V_{2} = p_{3}V_{3} =\ldots
		}
	\end{center}
\end{itemize}}{includegraphics[width=4cm]{Images/BoyleLaw}}
\subsubsection{Đường đẳng nhiệt}
immini[thm]{
	\begin{itemize}
		\item Đường biểu diễn sự biến thiên của áp suất theo thể tích khi nhiệt độ không đổi gọi là đường đẳng nhiệt.
		\item Dạng đường đẳng nhiệt:\\
		Trong hệ toạ độ (p, V) đường đẳng nhiệt là một nhánh của đường hyperbol.
		\begin{note}
			Trong các hệ trục tọa độ (V, T) và (p, T) thì đường đẳng nhiệt là đường thẳng vuông góc với trục T.
		\end{note}
	\end{itemize}
}{	\begin{tikzpicture}[scale=1,thick,>=stealth']
		\draw[->] (0,0)--(3.5,0)node[below]{$V$}; 
		\draw[->] (0,0)--(0,3)node[right]{p};
		\draw [blue, thick] (0.75,2.5) to [bend right=30] (2.8,0.75)node[right,black]{$T_2$};
		\draw [blue, thick] (0.25,2.5) to [bend right=35] (2.8,0.25)node[right,black]{$T_1$};
		\node at (2,2){$T_2>T_1$};
\end{tikzpicture}}
